\section{Proposition 4.4: failure at arbitrarily small angles}
\label{sec:prop44}

Write $P$ and $Q$ for the orthogonal projections onto two subspaces $U$ and
$V$. In the finite-dimensional acute case the canonical direct rotation is
the orthogonal factor in the polar decomposition of
\[
 T=QP+(I-Q)(I-P).
\]
A unitarily invariant norm depends only on singular values; in particular,
the trace norm is $\|X\|_1=\sum_j s_j(X)$. Proposition~4.4 of
\citet{davis1970rotation} asserts that, for real subspaces with largest
principal angle at most $\pi/3$, the direct rotation minimizes $\|I-W\|$
over orthogonal maps $W$ satisfying $WU=V$, for every such norm.
A real finite-dimensional counterexample therefore refutes the printed claim.

\subsection{An explicit one-parameter family}

\begin{theorem}[Trace-norm failure throughout the acute angle range]
\label{thm:small-angle}
For every $0<\theta<\pi/2$, there are two-dimensional subspaces
$U,V_\theta\subset\mathbb R^4$ and an orthogonal map $W_\theta$ with
$W_\theta U=V_\theta$, such that both principal angles equal $\theta$ and
the canonical direct rotation $R_\theta$ satisfies
\begin{align}
 \|I-W_\theta\|_1&=4\sin\theta,\label{eq:competitor-trace}\\
 \|I-R_\theta\|_1&=8\sin(\theta/2)>4\sin\theta.
 \label{eq:direct-trace}
\end{align}
\end{theorem}

\begin{proof}
Put $c=\cos\theta$, $s=\sin\theta$, and let
\[
 J=\begin{pmatrix}0&-1\\1&0\end{pmatrix},\quad
 C=\cos(2\theta)I_2+\sin(2\theta)J,\quad
 H=\frac{1}{\sqrt2}\begin{pmatrix}I_2&I_2\\I_2&-I_2\end{pmatrix}.
\]
Define
\begin{equation}
 W_\theta=H\begin{pmatrix}C&0\\0&I_2\end{pmatrix}H^T
 =\begin{pmatrix}c^2I_2+scJ&-s^2I_2+scJ\\
                  -s^2I_2+scJ&c^2I_2+scJ\end{pmatrix}.
 \label{eq:wtheta}
\end{equation}
Both $H$ and $C$ are orthogonal, so $W_\theta$ is orthogonal. Let
$U=\mathbb R^2\oplus\{0\}$ and $V_\theta=W_\theta U$.
The first two columns of $W_\theta$ factor as
\[
 \begin{pmatrix}c^2I_2+scJ\\-s^2I_2+scJ\end{pmatrix}
 =\begin{pmatrix}cI_2\\sJ\end{pmatrix}(cI_2+sJ).
\]
Since $cI_2+sJ$ is orthogonal, the columns of
$Y=\binom{cI_2}{sJ}$ form an orthonormal basis of $V_\theta$.
With $X=\binom{I_2}{0}$, one has $X^TY=cI_2$.
Thus both principal cosines are $c>0$, and both principal angles are
$\theta$.

The corresponding projections and canonical intertwiner are
\[
 P=\begin{pmatrix}I_2&0\\0&0\end{pmatrix},\qquad
 Q=YY^T=\begin{pmatrix}c^2I_2&-scJ\\scJ&s^2I_2\end{pmatrix},
\]
\begin{equation}
 T=QP+(I-Q)(I-P)
   =c\begin{pmatrix}cI_2&sJ\\sJ&cI_2\end{pmatrix}.
 \label{eq:intertwiner-theta}
\end{equation}
The last block matrix is orthogonal and $c>0$, so $|T|=cI_4$ and
its polar factor is exactly
\[
 R_\theta=\begin{pmatrix}cI_2&sJ\\sJ&cI_2\end{pmatrix}.
\]
In particular, this calculation identifies the canonical direct rotation
of the constructed pair. Since $R_\theta+R_\theta^T=2cI_4$,
\[
 (I-R_\theta)^T(I-R_\theta)=(2-2c)I_4.
\]
Consequently all four displacement singular values equal
$\sqrt{2-2c}=2\sin(\theta/2)$.

On the other hand, \cref{eq:wtheta} shows that the singular values of
$I-W_\theta$ are those of $(I_2-C)\oplus0_2$.
The identity
$(I_2-C)^T(I_2-C)=4s^2I_2$ gives the list
$(2s,2s,0,0)$. The two trace formulas follow. Finally,
\[
 4\sin\theta
 =8\sin(\theta/2)\cos(\theta/2)
 <8\sin(\theta/2),
\]
because $0<\theta/2<\pi/4$.
\end{proof}

\begin{corollary}[No positive short-angle threshold]
\label{cor:no-threshold}
For every $\varepsilon>0$, there are real two-dimensional subspaces of
$\mathbb R^4$ whose largest principal angle is strictly between zero and
$\varepsilon$ and is at most $\pi/3$, together with an orthogonal competitor
having strictly smaller trace-norm displacement than the direct rotation.
\end{corollary}

\begin{proof}
Apply \cref{thm:small-angle} with
$\theta=\min\{\varepsilon/2,\pi/4\}$.
\end{proof}

The direct rotation distributes its rotation over two principal planes;
$W_\theta$ rotates the diagonal plane $\{(x,x):x\in\mathbb R^2\}$ by
$2\theta$ and fixes its orthogonal complement. The latter displacement has
rank two, whereas the direct displacement has rank four. For this pair,
\[
 \frac{\|I-W_\theta\|_1}{\|I-R_\theta\|_1}=\cos(\theta/2).
\]
The ratio tends to one as $\theta$ tends to zero, but the inequality is strict
at every positive angle in the construction. The competitor is used only to
refute minimality; no claim that it is a global trace-norm minimizer is needed.
The operator norm exhibits the opposite comparison:
$\|I-W_\theta\|_{\rm op}=2\sin\theta$
exceeds $\|I-R_\theta\|_{\rm op}=2\sin(\theta/2)$.
The failure therefore depends on the norm class.

\subsection{The checked short-angle witness and the failed proof step}

At $\theta=\pi/4$, the family gives
\[
 W_{\pi/4}=\frac12
 \begin{pmatrix}
 1&-1&-1&-1\\1&1&1&-1\\-1&-1&1&-1\\1&-1&1&1
 \end{pmatrix}.
\]
The two trace norms are
\[
 \|I-W_{\pi/4}\|_1=2\sqrt2
 <4\sqrt{2-\sqrt2}=\|I-R_{\pi/4}\|_1.
\]
The construction, acuteness, principal-angle bound, singular-value
calculations, and strict inequality for this witness are checked in Lean.
Appendix~\ref{app:small-angle-artifact} identifies the corresponding files
and the separate parameterized proof source.

For $K=I-W_{\pi/4}$, let $\Omega_1$ and $\Omega_2$ be the projections onto
$\operatorname{span}\{e_0,e_3\}$ and
$\operatorname{span}\{e_1,e_2\}$. These are principal planes for the pair.
Davis and Kahan's equation~(4.3), applied to this configuration, requires
\begin{equation}
 \|K\|_{(4)}\geq
 \|K\Omega_1\|_{(2)}+\|K\Omega_2\|_{(2)},
 \label{eq:dk43-false}
\end{equation}
where $\|\cdot\|_{(j)}$ denotes the Ky Fan $j$-norm.
The checked calculations give
\[
 \|K\|_{(4)}=2\sqrt2,\qquad
 \|K\Omega_1\|_{(2)}=\|K\Omega_2\|_{(2)}=2.
\]
Thus \cref{eq:dk43-false} would imply $2\sqrt2\geq4$.
The domain planes are orthogonal, but their images under $K$ need not be;
the singular-value sums of these restrictions cannot be added in the
asserted direction.

\subsection{What the squared-displacement theorem does imply}

Call a unitarily invariant norm $N$ a $Q$-norm when there is a unitarily
invariant norm $M$ satisfying
\begin{equation}
 N(Z)^2=M(Z^*Z)\qquad\text{for every }Z.
 \label{eq:qnorm-def}
\end{equation}

\begin{proposition}[$Q$-norm consequence of Davis--Kahan Proposition~4.3]
\label{prop:qnorm-repair}
Let $U,V$ be an acute pair of subspaces of a finite-dimensional real or
complex Hilbert space, and let $R$ be their canonical direct rotation.
For every $Q$-norm $N$ and every orthogonal or unitary map $W$ with $WU=V$,
\begin{equation}
 N(I-R)\leq N(I-W).
 \label{eq:qnorm-repair}
\end{equation}
No $\pi/3$ restriction is required.
\end{proposition}

\begin{proof}
The valid squared-displacement theorem gives
\[
 M((I-R)^*(I-R))\leq M((I-W)^*(I-W)).
\]
By \cref{eq:qnorm-def}, these quantities are $N(I-R)^2$ and $N(I-W)^2$.
Both norm values are nonnegative, so taking square roots proves the claim.
\end{proof}

This implication is formalized for the finite-dimensional real/complex
unitarily invariant seminorm interface. The trace norm on the counterexample
space is also formally shown to fail the $Q$-norm predicate. The formal
result makes no identification of that predicate with additional named norm
families unless the corresponding identity is proved.

\subsection{Relationship to prior norm-minimization results}

Bolotnikov, Li, and Rodman \citep[Theorem~3.1]{bolotnikov2001minimal}
study unitary completion and prove a $Q$-norm minimization principle.
Their Lemma~3.3 gives a $2\times2$ example with a right-angle rotation and
complex diagonal competitors showing failure for Schatten $p$-norms with
$1\leq p<2$. These results already establish the need to distinguish
norm classes. The $Q$-norm principle therefore precedes the present
formalization. The right-angle example with complex diagonal competitors
does not itself contradict the real short-angle restriction.

Qiu, Zhang, and Li \citep[Section~5]{qiu2006grassmann} list three extremal
properties attributed to Davis and Kahan. Their third item asserts
full-displacement minimality under the $\pi/3$ bound for every unitarily
invariant norm. The real family in \cref{thm:small-angle} contradicts that
item as well as the original proposition. Its conclusion concerns this
specific extremal assertion; it does not refute the sine-theta bounds or
assert a failure of the main Grassmann-metric results in that article.

Zhang and Qiu \citep[Section~4, p.~1615]{zhang2010canonicalangles}
state that direct rotations minimize Frobenius-norm displacement among
unitary transformations mapping one subspace onto another. They explicitly
note that direct rotations need not minimize trace-norm displacement and
that spectral-norm displacement can have additional minimizers. Their
statement about trace-norm failure is general and does not address the
specific real short-angle hypothesis of Davis--Kahan Proposition~4.4.

The real family in \cref{thm:small-angle} satisfies every positive
angle threshold, and \cref{eq:dk43-false} identifies the failed inequality
in the original argument. The $Q$-norm consequence
\cref{prop:qnorm-repair} follows from Davis and Kahan's valid
squared-displacement theorem and is consistent with the earlier norm-class
results of \citet{bolotnikov2001minimal}.
